\color{unidarkgreen}{\subsection{Защита при потере охлаждения}}
\color{black}

\subsubsection{Общие сведения}

\begin{enumerate}[label=\arabic{section}.\arabic{subsection}.\arabic{subsubsection}.\arabic*,  align=left, labelsep=4pt, leftmargin=0pt, itemindent=57pt]

\item
Защита при потере охлаждения направлена на предотвращение перегрева и повреждения
трансформатора в результате частичного или полного отказа системы охлаждения.

\item
Функция защиты при потере охлаждения контролирует уровни тока на обмотках ВН, НН1(НН),
НН2(СН). К каждой из обмоток подключено по 2 реле тока для контроля за уровнем нагрузки
автотрансформатора. Каждое реле имеет индивидуально настраиваемый порог по току срабатывания,
задаваемый пользователем.

\item
Защита содержит 4 ступени, каждая из которых может быть выполнена с контролем температуры
масла и уровня тока. Ввод контроля перегрева масла, а также выбор соответствующего уровня тока
осуществляется программными накладками \textbf{«ЗПО-N конт.ДТм»} и \textbf{«ЗПО-N конт.РТ»} соответственно
(где N – номер ступени). Каждая из ступеней имеет индивидуальную регулируемую задержку на
срабатывание.

\item
На рисунке \ref{pic:ZPO} представлена защита при потере охлаждения 1 ступень, последующие аналогичны. 

\begin{figure}[H]
\centering
\includegraphics[width=0.7\textwidth,height=0.7\textheight,keepaspectratio]{\fbpath/011.1161 ЗПО/img/ZPO_1.pdf}
%\captionof{figure}{Функциональный блок ЗПО}\label{pic:ZPO}
\end{figure}

\begin{figure}[H]
\centering
\includegraphics[width=0.7\textwidth,height=0.7\textheight,keepaspectratio]{\fbpath/011.1161 ЗПО/img/ZPO_2.pdf}
%\captionof{figure}{Функциональный блок ЗПО}\label{pic:ZPO}
\end{figure}

\begin{figure}[H]
\centering
\includegraphics[width=0.7\textwidth,height=0.7\textheight,keepaspectratio]{\fbpath/011.1161 ЗПО/img/ZPO_3.pdf}
\captionof{figure}{Функциональный блок ЗПО}\label{pic:ZPO}
\end{figure}


\small
\begin{longtable}{|>{\centering\arraybackslash}m{5.3cm}|>{\centering\arraybackslash}m{4.2cm}|>{\centering\arraybackslash}m{1cm}|>{\centering\arraybackslash}m{3.3cm}|>{\centering\arraybackslash}m{1.8cm}|}
% Первая страница: полный заголовок
\caption{Перечень уставочных параметров функции пуска охладителей\hfill\vspace{-0.5\baselineskip}}\label{desc:tbl1}\\
\hline
\rowcolor{gray!30}
Наименование уставки & Диапазон & Шаг & Значение по умолчанию & Единица измерения \\
\hline
\endfirsthead

% Промежуточные страницы: «Продолжение таблицы»
\caption*{\hspace{3pt}\emph{Продолжение таблицы \ref{desc:tbl1}}\hfill\ }\vspace{-0.5\baselineskip} \\
\hline
\rowcolor{gray!30}
Наименование уставки & Диапазон & Шаг & Значение по умолчанию & Единица измерения \\
\hline
\endhead

%===m>ДТЗ|-
%===m
\end{longtable}
\normalsize	

\end{enumerate}